\documentclass[a4paper,12pt]{article}
\usepackage[T2A]{fontenc}
\usepackage[utf8]{inputenc}
\usepackage[russian]{babel}
\usepackage{amsmath,amssymb,mathtools}
\usepackage{geometry}
\geometry{left=23mm,right=23mm,top=24mm,bottom=22mm}
\usepackage{microtype}
\usepackage{icomma}
\usepackage{setspace}
\setstretch{1.16}
\usepackage{parskip}
\usepackage{enumitem}
\setlist[itemize]{leftmargin=1.5em,itemsep=.27em,topsep=.3em}
\setlist[enumerate]{leftmargin=1.6em,itemsep=.55em,topsep=.35em}
\usepackage{tabularx,booktabs,array}
\usepackage{xcolor}
\definecolor{accent}{HTML}{087A7E}
\definecolor{accentdark}{HTML}{0B5057}
\definecolor{ink}{HTML}{292D30}
\definecolor{muted}{HTML}{5E666A}
\color{ink}
\usepackage{tikz}
\usetikzlibrary{arrows.meta,positioning,shapes.geometric,calc}
\usepackage{fancyhdr}
\usepackage[hidelinks]{hyperref}
\hypersetup{pdftitle={Квадратные корни: чек-лист},pdfauthor={Лёвин Артём Александрович},pdfsubject={Алгебра, 8 класс, Екатерина Скелли}}
\pagestyle{fancy}
\fancyhf{}
\renewcommand{\headrulewidth}{0pt}
\fancyhead[L]{\footnotesize\textcolor{muted}{Квадратные корни}}
\fancyhead[R]{\footnotesize\textcolor{muted}{Екатерина Скелли}}
\fancyfoot[C]{\footnotesize\textcolor{muted}{\thepage}}
\setlength{\headheight}{15pt}
\newcommand{\key}[1]{\par\textcolor{accentdark}{\textbf{Главная идея.}}~#1\par}
\newcommand{\warn}[1]{\par\textcolor{accentdark}{\textbf{Проверьте.}}~#1\par}
\newcommand{\example}[1]{\noindent\textcolor{accentdark}{\textbf{Пример.}}~#1}
\newcommand{\ruleline}{\par\vspace{.3em}\noindent\textcolor{accent}{\rule{\linewidth}{0.55pt}}\par\vspace{.35em}}
\setcounter{secnumdepth}{1}
\renewcommand{\thesection}{\arabic{section}}
\makeatletter
\renewcommand{\section}{\@startsection{section}{1}{\z@}{2.1ex plus .8ex minus .4ex}{.9ex}{\normalfont\Large\bfseries\color{accentdark}}}
\makeatother
\begin{document}
\thispagestyle{empty}
\begin{center}
\vspace*{2.7cm}
{\large\textcolor{accentdark}{МАТЕМАТИКА}}\par
\vspace{1.5cm}
{\LARGE Чек-лист}\par
\vspace{1.05cm}
{\LARGE\textcolor{accent}{Квадратные корни:}}\par
\vspace{.12cm}
{\LARGE\textcolor{accent}{свойства и преобразования}}\par
\vspace{1.4cm}
{\large Алгебра · 8 класс}\par
\vfill
{\large Екатерина Скелли}\par
\vspace{.8cm}
{20 ноября 2026 г.}\par
\vspace{1.2cm}
{\small Лёвин Артём Александрович}\par
{\small эксклюзивно для Екатерины Скелли}\par
\vspace{1.2cm}
\begin{tikzpicture}[x=1cm,y=1cm]
\draw[accent,line width=.7pt] (0,0) .. controls (1.25,.47) and (2.6,-.37) .. (4.1,0)
.. controls (5.3,.42) and (6.55,-.43) .. (8.2,0);
\fill[accent] (4.1,0) circle (1.1pt);
\end{tikzpicture}
\end{center}
\clearpage

\section{Что необходимо помнить о корне}
\key{Квадратный корень \(\sqrt a\) --- это неотрицательное число, квадрат которого равен \(a\).}
\[
\sqrt a\ge0,\qquad (\sqrt a)^2=a\quad(a\ge0),\qquad
\sqrt{a^2}=|a|.
\]
В частности, \(\sqrt{49}=7\), \(\sqrt{0}=0\), \(\sqrt{2}\) --- точное значение, его необязательно превращать в десятичную дробь.

\warn{В действительных числах для \(\sqrt a\) требуется \(a\ge0\). Если корень стоит в знаменателе, его значение должно быть строго положительным.}

\section{Умножение и деление квадратных корней}
При \(a\ge0\), \(b\ge0\):
\[
\boxed{\sqrt a\,\sqrt b=\sqrt{ab}}.
\]
При \(a\ge0\), \(b>0\):
\[
\boxed{\frac{\sqrt a}{\sqrt b}=\sqrt{\frac ab}}.
\]
Удобно сначала объединить корни, затем сократить дробь под корнем.

\example{Вычислите \(\sqrt{\frac23}\cdot\sqrt{\frac{27}{8}}\).}
\[
\sqrt{\frac23\cdot\frac{27}{8}}
=\sqrt{\frac94}=\frac32=1{,}5.
\]
\example{Вычислите \(\displaystyle\frac{\sqrt{26}}{\sqrt{14}\sqrt{91}}\).}
\[
\sqrt{\frac{26}{14\cdot91}}
=\sqrt{\frac{2\cdot13}{(2\cdot7)(7\cdot13)}}
=\sqrt{\frac1{49}}=\frac17.
\]
\warn{Для суммы корней такое объединение неверно: \(\sqrt a+\sqrt b\ne\sqrt{a+b}\) в общем случае.}

\section{Смешанные числа под знаком корня}
Перед вычислением произведения смешанные числа переводят в неправильные дроби:
\[
\begin{aligned}
&\sqrt{\left(2\frac{23}{49}\right)\left(2\frac{2}{49}\right)}\\
&=\sqrt{\frac{121}{49}\cdot\frac{100}{49}}
=\frac{11\cdot10}{7\cdot7}=\frac{110}{49}.
\end{aligned}
\]
\key{Сначала преобразуйте смешанные числа, затем умножайте и извлекайте корень.}

\section{Вынесение множителя из-под знака корня}
Если под корнем есть полный квадрат, его можно вынести:
\[
\boxed{\sqrt{m^2b}=|m|\sqrt b\qquad(b\ge0)}.
\]
Для положительного \(m\) модуль можно опустить.
\[
\sqrt{20}=\sqrt{4\cdot5}=2\sqrt5,\qquad
\sqrt{500}=\sqrt{100\cdot5}=10\sqrt5.
\]
\[
\sqrt{18}=\sqrt{3^2\cdot2}=3\sqrt2,\qquad
\sqrt{162}=\sqrt{9^2\cdot2}=9\sqrt2.
\]
\warn{Корень извлекается из полного квадрата; например, \(\sqrt{3^2\cdot2}=3\sqrt2\).}

\section{Внесение множителя под знак корня}
Для \(c\ge0\) и \(a\ge0\):
\[
\boxed{c\sqrt a=\sqrt{c^2a}}.
\]
Внутри корня множитель появляется \emph{в квадрате}:
\[
2\sqrt8=\sqrt{2^2\cdot8}=\sqrt{32},\qquad
4\sqrt7=\sqrt{16\cdot7}=\sqrt{112}.
\]
\warn{Если \(c<0\), знак минуса должен остаться перед корнем: \(-3\sqrt2=-\sqrt{18}\).}

\section{Подобные корни и упрощение выражений}
\key{Складывать и вычитать коэффициенты можно только при одинаковых подкоренных выражениях.}
\[
\begin{aligned}
\sqrt{20}-\sqrt{500}+\sqrt{75}
&=2\sqrt5-10\sqrt5+5\sqrt3\\
&=(2-10)\sqrt5+5\sqrt3\\
&=\boxed{-8\sqrt5+5\sqrt3}.
\end{aligned}
\]
Общий корень можно вынести за скобки:
\[
2\sqrt5-10\sqrt5=(2-10)\sqrt5=-8\sqrt5.
\]
Ещё один приём --- разложение корня для сокращения дроби:
\[
\begin{aligned}
\frac{\sqrt{14}-\sqrt{28}}{\sqrt7}
&=\frac{\sqrt7(\sqrt2-2)}{\sqrt7}
=\boxed{\sqrt2-2}.
\end{aligned}
\]
Здесь \(\sqrt7\ne0\), поэтому сокращение допустимо.

\section{Квадрат корня и разность квадратов}
\[
(\sqrt a)^2=a\quad(a\ge0),\qquad
(u-v)(u+v)=u^2-v^2.
\]
\example{Вычислите \((\sqrt{17}-\sqrt3)(\sqrt{17}+\sqrt3)\).}
\[
(\sqrt{17})^2-(\sqrt3)^2=17-3=\boxed{14}.
\]
\example{Вычислите \(\sqrt{7^4\cdot2^8}\).}
\[
\sqrt{(7^2)^2\cdot(2^4)^2}=7^2\cdot2^4=49\cdot16=\boxed{784}.
\]
\warn{В \((2\sqrt3)^2\) квадрат относится ко всему произведению: \(4\cdot3=12\).}

\section{Сравнение выражений с корнями}
Если оба сравниваемых выражения неотрицательны, их можно возвести в квадрат. Знак сравнения сохранится.
\[
7=\sqrt{49},\qquad 5\sqrt2=\sqrt{25\cdot2}=\sqrt{50}.
\]
Поэтому \(49<50\), значит
\[
\boxed{7<5\sqrt2}.
\]
Аналогично \((3\sqrt2)^2=18\), \((2\sqrt3)^2=12\), значит
\[
\boxed{3\sqrt2>2\sqrt3}.
\]
\key{У сравниваемых неотрицательных чисел большему квадрату соответствует большее число.}

\section{Умножение: множители и скобки}
Различайте произведение и сумму внутри скобок:
\[
2(2\sqrt3)=4\sqrt3,
\qquad
2(2+\sqrt3)=4+2\sqrt3.
\]
При умножении на сумму умножаем на \emph{каждое} слагаемое:
\[
\begin{aligned}
(4\sqrt3)(5+\sqrt3)
&=20\sqrt3+4(\sqrt3)^2\\
&=\boxed{20\sqrt3+12}.
\end{aligned}
\]
При умножении двух сумм получаются четыре произведения:
\[
\begin{aligned}
(4+\sqrt3)(5+\sqrt3)
&=4\cdot5+4\sqrt3+5\sqrt3+(\sqrt3)^2\\
&=20+9\sqrt3+3\\
&=\boxed{23+9\sqrt3}.
\end{aligned}
\]
\warn{Если в скобках стоит плюс, применяется распределительный закон. Если стоит только умножение, можно переставить и сгруппировать множители.}

\section{Корень в знаменателе дроби}
Освободиться от иррациональности в знаменателе --- значит преобразовать дробь так, чтобы в знаменателе не было корня.
\[
\frac1{\sqrt7}=\frac{1\cdot\sqrt7}{\sqrt7\cdot\sqrt7}
=\boxed{\frac{\sqrt7}{7}}.
\]
Числитель и знаменатель домножаются на одно и то же \emph{ненулевое} число, поэтому дробь сохраняет значение.
\[
\frac5{\sqrt5}=\frac{5\sqrt5}{5}=\boxed{\sqrt5},
\qquad
\frac3{\sqrt3}=\boxed{\sqrt3}.
\]
\example{Квадрат произведения в знаменателе:}
\[
\frac6{(2\sqrt3)^2}=\frac6{2^2(\sqrt3)^2}
=\frac6{4\cdot3}=\boxed{\frac12}.
\]
\key{При дроби \(\frac a{\sqrt a}\), где \(a>0\), результат равен \(\sqrt a\).}

\clearpage
\thispagestyle{fancy}
\begin{center}{\Large\bfseries\textcolor{accentdark}{Алгоритм: как упростить выражение с корнями}}\end{center}
\vspace{1.0cm}
\begin{center}
\begin{tikzpicture}[
node distance=9mm,
start/.style={ellipse,draw=accentdark,thick,minimum width=39mm,minimum height=10mm,align=center,fill=white},
action/.style={rectangle,rounded corners=2pt,draw=accentdark,thick,text width=111mm,minimum height=12mm,align=center,inner sep=6pt,fill=white},
finish/.style={ellipse,draw=accentdark,thick,minimum width=45mm,minimum height=10mm,align=center},
arr/.style={-{Stealth[length=3mm]},line width=.85pt,accentdark}]
\node[start] (a) {Начало};
\node[action,below=of a] (b) {Проверьте условия: подкоренные выражения \(\ge 0\), знаменатели \(\ne0\)};
\node[action,below=of b] (c) {Найдите полные квадраты под корнями и вынесите их};
\node[action,below=of c] (d) {Если есть произведение или дробь корней, используйте свойства корня};
\node[action,below=of d] (e) {Если есть скобки, примените подходящую формулу или распределительный закон};
\node[action,below=of e] (f) {Приведите подобные корни с одинаковыми подкоренными выражениями};
\node[action,below=of f] (g) {Уберите корень из знаменателя, если он там остался};
\node[finish,below=of g] (h) {Проверьте результат};
\draw[arr] (a)--(b);\draw[arr] (b)--(c);\draw[arr] (c)--(d);
\draw[arr] (d)--(e);\draw[arr] (e)--(f);\draw[arr] (f)--(g);\draw[arr] (g)--(h);
\end{tikzpicture}
\end{center}
\clearpage

\section{Частые ошибки и мини-памятка}
\begin{tabularx}{\linewidth}{@{}>{\raggedright\arraybackslash}p{.36\linewidth}X@{}}
\toprule
\textbf{Ошибочное действие} & \textbf{Правильный ориентир}\\
\midrule
\(\sqrt{a+b}=\sqrt a+\sqrt b\) & Для суммы такого свойства нет. Сначала вычислите сумму под корнем.\\[.45em]
\(\sqrt{a^2}=a\) всегда & Для действительного \(a\): \(\sqrt{a^2}=|a|\).\\[.45em]
\(2\sqrt3+3\sqrt5=5\sqrt8\) & Разные подкоренные выражения не складываются.\\[.45em]
\((2\sqrt3)^2=2\cdot3\) & Возведите в квадрат оба множителя: \(4\cdot3=12\).\\[.45em]
\(2(3+\sqrt5)=6+\sqrt5\) & Двойка умножается на оба слагаемых: \(6+2\sqrt5\).\\[.45em]
\(\frac1{\sqrt5}\cdot\frac{\sqrt2}{\sqrt2}\) & Корень в знаменателе сохранится; домножайте на \(\sqrt5/\sqrt5\).\\
\bottomrule
\end{tabularx}
\vspace{.9cm}
\textcolor{accentdark}{\textbf{Перед завершением решения проверьте:}}
\begin{itemize}
\item все ли полные квадраты вынесены из-под корней;
\item можно ли привести подобные корни;
\item правильно ли учтены скобки, коэффициенты и знаки;
\item нет ли корня в знаменателе;
\item соблюдены ли допустимые значения выражений.
\end{itemize}

\clearpage
\begin{center}{\Large\bfseries\textcolor{accentdark}{Тренировочный блок}}\end{center}
\vspace{.3cm}
\noindent\textcolor{muted}{10 заданий · от простого к сложному. Записывайте основные преобразования.}
\ruleline
\begin{enumerate}
\item Вычислите: \(\displaystyle\sqrt{\frac35}\cdot\sqrt{\frac{20}{27}}\).
\item Вычислите: \(\displaystyle\frac{\sqrt{21}}{\sqrt7\cdot\sqrt{12}}\).
\item Вычислите: \(\displaystyle\sqrt{\left(1\frac9{16}\right)\left(2\frac14\right)}\).
\item Упростите: \(\displaystyle\sqrt{72}-\sqrt{98}+\sqrt{50}\).
\item Вынесите множитель из-под знака корня: \(\sqrt{288}\).
\item Внесите множитель под знак корня: \(3\sqrt{11}\).
\item Вычислите: \((\sqrt{23}-\sqrt8)(\sqrt{23}+\sqrt8)\).
\item Сравните числа \(8\) и \(3\sqrt7\), поставив нужный знак.
\item Упростите: \((4+\sqrt3)(5+\sqrt3)\).
\item Освободитесь от иррациональности в знаменателях и упростите:
\[\frac6{\sqrt3}-\frac5{\sqrt5}.\]
\end{enumerate}
\vspace{.2cm}
\warn{Если получилось извлечь квадратный корень из целого полного квадрата, обязательно сделайте это.}

\clearpage
\begin{center}{\Large\bfseries\textcolor{accentdark}{Ответы}}\end{center}
\vspace{.4cm}
\begin{table}[h]
\centering
\renewcommand{\arraystretch}{1.5}
\begin{tabularx}{.85\linewidth}{@{}cX@{}}
\toprule
\textbf{№} & \textbf{Ответ}\\
\midrule
1 & \(\frac23\)\\
2 & \(\frac12\)\\
3 & \(\frac{15}{8}\)\\
4 & \(4\sqrt2\)\\
5 & \(12\sqrt2\)\\
6 & \(\sqrt{99}\)\\
7 & \(15\)\\
8 & \(8>3\sqrt7\)\\
9 & \(23+9\sqrt3\)\\
10 & \(2\sqrt3-\sqrt5\)\\
\bottomrule
\end{tabularx}
\end{table}
\vspace{.9cm}
\noindent\textcolor{accentdark}{\textbf{Итоговая памятка}}\par
\begin{tabularx}{\linewidth}{@{}>{\raggedright\arraybackslash}p{.39\linewidth}X@{}}
\toprule
\textbf{Сигнал в задаче} & \textbf{Полезное действие}\\
\midrule
Произведение корней & Объедините их под одним знаком корня.\\
Полный квадрат под корнем & Вынесите его множитель наружу.\\
Сумма подобных корней & Посчитайте коэффициенты.\\
Сопряжённые скобки & Используйте разность квадратов.\\
Сравнение неотрицательных чисел & Сравните квадраты.\\
Корень в знаменателе & Домножьте дробь на подходящий корень.\\
\bottomrule
\end{tabularx}
\vfill
\begin{center}\footnotesize\textcolor{muted}{Лёвин Артём Александрович · эксклюзивно для Екатерины Скелли}\end{center}
\end{document}
